\chapter{RESULTADOS ESPERADOS}

Com base na metodologia proposta, espera-se o desenvolvimento de um protótipo funcional de sistema semiautomático baseado em \gls{ia} para a validação de planos de ensino no curso de Ciência da Computação do \gls{ifc} Videira. Os principais resultados incluem:
\begin{itemize}
\item um mapeamento completo dos elementos obrigatórios dos planos de ensino, alinhados à Organização Didática do \gls{ifc} e ao \gls{ppc}, estruturado como base de conhecimento normativa consultável pelo sistema;
\item um \textit{pipeline} que detecte lacunas e gere sugestões de correção por critério, com \textit{recall} e \textit{precision} satisfatórios em testes comparativos com a validação manual do \gls{nde};
\item um painel web que permita acompanhar os planos processados e o resultado de cada critério, reduzindo a sobrecarga administrativa;
\item uma avaliação dos impactos administrativos e pedagógicos do sistema, incluindo o retorno qualitativo de docentes e membros do \gls{nde}.
\end{itemize}

Os resultados contribuirão para a inovação educacional, promovendo automação escalável e alinhada às normas institucionais, com potencial para expansão a outros cursos do \gls{ifc}.
